\chapter{Experimentation and Evaluation}
\label{ch:results}

This chapter evaluates the score formula, the divergence-point detector, and the use of their feedback in complete refactoring sessions. The formula is tested through sensitivity and ablation experiments, the detector against labelled injection sessions, and the end-to-end workflow through a 30-session human study plus an agent extension on the same codebase.

\S\ref{sec:results-metric} reports the score-formula experiments. \S\ref{sec:results-detector} reports detector precision and recall against externally labelled divergence kinds. \S\ref{sec:results-studies} reports the human and agent studies, where the outcome is whether feedback is reflected in later refactoring behaviour.

The labelled detector experiment uses the 45-session injection set and two-label scheme defined in \S\ref{sec:methodology-rater}. Because the per-kind counts are small, tables report raw counts alongside percentages. Every magnitude value in this chapter is the trajectory-final $J$ delta defined in \S\ref{sec:methodology-score}, so magnitudes are comparable across the four divergence kinds.

\section{Testing the score formula}
\label{sec:results-metric}

The score weights were chosen from a literature-supported severity ordering (\S\ref{sec:methodology-weights}), not fitted to an external outcome dataset. This section therefore tests whether the rankings are stable under plausible weight changes and whether the seven process-side terms contribute measurable signal, following the evaluation design in \S\ref{sec:methodology-score-eval}. \S\ref{sec:results-headline} then reports a structural limit of the formula that appears in both experiments.

\subsection{Sensitivity sweep}
\label{sec:results-sensitivity}

The first experiment measures how sensitive the divergence-point ranking is to the weights in the process-score formula. Large ranking changes under small weight perturbations would indicate that the reported priorities depend too strongly on arbitrary parameter choices. The experiment uses the single-knob and multi-knob perturbation methods defined in \S\ref{sec:methodology-sensitivity-design}. The main result is that the top-ranked divergence point is usually stable across all three session sets, although some terms do affect the ranking in identifiable cases.

\paragraph{Experiment setup.}
For the single-knob sweep, the chosen weight is scaled by one of nine factors in $\{0.1, 0.25, 0.5, 0.75, 1.25, 1.5, 2.0, 4.0, 10.0\}$. With twelve weights ($6$ process-side and $6$ cleanliness-side), this gives $12 \times 9 \times \lvert \mathcal{C} \rvert$ perturbation runs per session set $\mathcal{C}$. The multi-knob sweep draws $200$ independent samples per session using the log-normal sampling scheme defined in \S\ref{sec:methodology-sensitivity-design}.

The experiment is run against three session sets: the $45$ hand-labelled injection sessions of \S\ref{sec:results-divergence}, the $30$ user-study sessions of \S\ref{sec:results-userstudy}, and the $48$ agent extension sessions of \S\ref{sec:results-userstudy-agent}. These differ in their per-session divergence-point profiles. The injection sessions have at most $5$ divergence points per session and concentrate on single-kind exercises. The user-study sessions have up to $13$ per session across naturalistic sessions ($121$ in total, mean $4.0$ per session, with the no-feedback baseline arm carrying most of the high-count sessions). The agent sessions have at most $3$ per session across the eight agent stacks ($42$ in total, mean $0.9$ per session), consistent with the structural detector biases against agent edit patterns discussed in \S\ref{sec:results-userstudy-agent}.

As defined in \S\ref{sec:methodology-sensitivity-design}, ranking stability is reported using Kendall's $\tau$-b and top-$N$ hit rate. Top-$1$ is the only metric that works across all three session sets, since it applies to any session with at least one divergence point. Top-$3$ and top-$5$ are reported on the user-study sessions only, where the per-session rankings are large enough to make them informative. Several divergence-point kinds produce ties, so ties are broken by ascending step index in both the production and perturbed rankings.

\paragraph{Single-knob results.}
Table~\ref{tab:results-sensitivity-headline} reports top-$1$ stability across the three session sets. The user-study sessions are the most demanding test: they have the longest per-session rankings, the highest divergence-point density, and the lowest baseline saturation. The injection sessions show the highest top-$1$ stability, but only because the per-session rankings are shorter and offer less surface area to reshuffle. The agent sessions sit between them, and $17.5\%$ of their baseline divergence points hit the score-clamp ceiling, which suppresses the observable effect of weight perturbations on a non-trivial slice of the agent set. On the user-study sessions the top-$1$ recommendation is preserved on $3{,}141$ of $3{,}240$ rows ($96.9\%$).

\begin{table}[htbp]
\centering
\small
\begin{tabular}{lrrrrr}
\hline
\textbf{session set} & \textbf{rows} & \textbf{$\tau = 1.0$} & \textbf{top-$1$ $= 1$} & \textbf{$\tau < 0$} & \textbf{baseline saturation} \\
\hline
Injection ($45$)         & $4{,}860$ & $77.0\%$ & $99.8\%$ & $0.3\%$ & $10.6\%$ \\
\textbf{User study ($30$)} & $3{,}240$ & $51.5\%$ & $\mathbf{96.9\%}$ & $0.8\%$ & $\mathbf{6.4\%}$ \\
Agent ($48$)              & $5{,}184$ & $88.5\%$ & $98.5\%$ & $1.2\%$ & $17.5\%$ \\
\hline
\end{tabular}
\caption{Single-knob sensitivity across three session sets. The user-study sessions are the discriminating benchmark: longest per-session rankings and lowest baseline saturation. The $\tau = 1.0$ share is meaningfully lower on the user-study set ($51.5\%$) than on the injection set ($77.0\%$) or the agent set ($88.5\%$) because user-study rankings are long enough to allow tail reshuffling under arbitrary single-knob perturbation, even when the top-of-ranking decision is preserved on $96.9\%$ of rows.}
\label{tab:results-sensitivity-headline}
\end{table}

\paragraph{Sources of ranking instability.}
Table~\ref{tab:results-sensitivity-knobs} reports per-knob top-$1$ disruption counts on the user-study sessions. Every process-side weight produces measurable disruption between $4.4\%$ and $8.9\%$ of rows; the six cleanliness sub-weights produce at most $0.7\%$ disruption each, with three of them producing none at all. The largest disruptor is the cleanliness-gain weight $W_{\textsc{g}}$ ($24$ of $270$ rows, $8.9\%$), followed by the manual-when-IDE weight $W_{\textsc{mi}}$ ($17$ rows), the skipped-tests weight $W_{\textsc{st}}$ ($15$), and the commit-gap weight $W_{\textsc{cg}}$ ($15$). The length-bonus and broken-build weights are close behind ($14$ and $12$). The decomposition tells us the score formula's \emph{effective parameter dimension} is closer to six (the process-side terms) than to twelve: the cleanliness composite usually adds the same amount to both paths so it does not change the ranking, and the process-side weights do almost all of the rank-shaping.

\begin{table}[htbp]
\centering
\small
\begin{tabular}{lrr}
\hline
\textbf{knob} & \textbf{top-$1$ disruptions / $270$ rows} & \textbf{share} \\
\hline
\texttt{gain}        & $24$ & $8.9\%$ \\
\texttt{manualIde}   & $17$ & $6.3\%$ \\
\texttt{skipTests}   & $15$ & $5.6\%$ \\
\texttt{commitGap}   & $15$ & $5.6\%$ \\
\texttt{length}      & $14$ & $5.2\%$ \\
\texttt{broken}      & $12$ & $4.4\%$ \\
\hline
\multicolumn{3}{l}{\textit{cleanliness sub-weights: \texttt{coupling} $2$ ($0.7\%$), \texttt{smells} $1$, \texttt{cohesion} $1$, all others $0$}} \\
\hline
\end{tabular}
\caption{Per-knob top-$1$ disruption counts on the user-study sessions (30 sessions $\times$ 9 perturbation factors $=$ 270 rows per knob). The six process-side weights account for almost every top-$1$ shift; cleanliness sub-weights produce at most $0.7\%$ disruption each. The cleanliness-gain weight $W_{\textsc{g}}$ is the most leverage-y single knob on this corpus.}
\label{tab:results-sensitivity-knobs}
\end{table}

\paragraph{Multi-knob Monte Carlo.}
The single-knob sweep only tests one weight at a time. The multi-knob Monte Carlo perturbs every weight together. Table~\ref{tab:results-sensitivity-mc} reports the distribution of $\tau$ and top-$N$ hit rate across $200 \times \lvert \mathcal{C} \rvert$ joint-perturbation samples per session set.

The multi-knob result is meaningfully weaker than the single-knob headline. On the user-study sessions, the top-$1$ recommendation changes on $13.1\%$ of joint-perturbation samples, against $3.1\%$ under single-knob, and mean $\tau$ falls to $0.656$. The injection sessions are more stable on top-$1$ ($99.1\%$) but their multi-knob mean $\tau$ falls to $0.805$, comparable in magnitude to the user-study drop, because joint perturbation can reshuffle long stretches of the tail even when the highest-magnitude divergence point is preserved. The agent sessions sit between the two on mean $\tau$ ($0.848$), with top-$1$ stability at $92.6\%$. The formula is robust to small joint perturbations and degrades smoothly with larger ones. It is not robust to arbitrary combinations of weights, even within the realistic range we tested.

\begin{table}[htbp]
\centering
\small
\begin{tabular}{lrrrrr}
\hline
\textbf{session set} & \textbf{samples} & \textbf{mean $\tau$} & \textbf{top-$1$ $= 1$} & \textbf{top-$3$ $= 1$} & \textbf{top-$5$ $= 1$} \\
\hline
Injection ($45$)         & $9{,}000$ & $0.805$ & $99.1\%$ & $77.8\%$ & $77.8\%$ \\
\textbf{User study ($30$)} & $6{,}000$ & $\mathbf{0.656}$ & $\mathbf{86.9\%}$ & $\mathbf{68.0\%}$ & $66.6\%$ \\
Agent ($48$)              & $9{,}600$ & $0.848$ & $92.6\%$ & $91.0\%$ & $89.6\%$ \\
\hline
\end{tabular}
\caption{Multi-knob Monte Carlo robustness across three session sets ($200$ samples per session, log-normal $\sigma = \ln 2$). The user-study sessions are the discriminating benchmark; the multi-knob top-$1$ stability is $\sim 10$ percentage points lower than the single-knob result on the same set. Top-$3$ and top-$5$ track top-$1$ on the agent set ($91.0\%$ and $89.6\%$) but drop substantially on the user-study set ($68.0\%$ and $66.6\%$) where per-session rankings are longer and the tail is more reshuffleable.}
\label{tab:results-sensitivity-mc}
\end{table}

\paragraph{Saturation and score-clamp bias.}
One concern is whether the $\tau$ and top-$N$ numbers reflect the $[0, 100]$ clamp on the process score rather than true stability. When both user and alt scores hit the ceiling, their magnitude is zero regardless of weights, which would inflate $\tau$. Saturation is therefore tracked per session. The right-hand column of Table~\ref{tab:results-sensitivity-headline} shows the three session sets differ substantially. The injection sessions have $10.6\%$ saturated baseline divergence points, concentrated on the short single-kind sessions. The agent sessions have $17.5\%$, because many agent sessions push the score close to the ceiling on the easier playbook sessions (the S4 no-op session in particular runs at $J = 45$ across every agent cell). The user-study sessions sit lowest at $6.4\%$. The user-study set is therefore the cleanest benchmark for the perturbation statistics: its $\tau$ and top-$N$ values are dominated by real perturbation response rather than clipping artefacts.

\paragraph{Caveats and scope.}
Three main limitations. First, the experiment establishes that the formula is stable near the production weights. It does not establish that the production weights themselves are uniquely correct, only that small changes to them mostly preserve the ranking. To properly validate this, we'd need to refit the weights against an external outcome dataset, e.g.\ longitudinal code-quality measurements, but no such dataset exists for this problem. Second, the multi-knob test is centred around the production weights and covers the neighbourhood within roughly $\pm 2\sigma$ (about $\times 0.25$ to $\times 4$). The formula visibly degrades on the user-study sessions even inside this range ($\tau = 0.656$ mean), so the robustness claim is local rather than global. Third, the user-study sessions are $n = 30$ across $5$ participants split into a with-feedback arm ($n = 3$) and a no-feedback baseline arm ($n = 2$); both arms feed the robustness corpus on the basis that they exercise the same scoring formula on the same codebase. We document these limits as threats to validity in \S\ref{sec:results-scoped-out} and the discussion chapter. The chapter does not claim more than the experiments support.

\subsection{Ablation sweep}
\label{sec:results-ablation}

The sensitivity sweep shows the ranking is robust to weight perturbations. This experiment asks the different question of whether every term in the score formula carries substantial weight. Using the ablation design defined in \S\ref{sec:methodology-ablation-design}, the experiment enumerates all $2^7 = 128$ subsets of the seven process-side weights (including the intermediate-cleanliness lag weight $W_{\textsc{lag}}$), producing $5{,}760$ rows across the $45$ injection sessions. The cleanliness sub-weights are kept at production because the sensitivity sweep showed that perturbing them almost never reshuffles the ranking; ablating them would inflate the variant count from $128$ to $8{,}192$ for little additional information. The sweep rules out two failure modes: sum-over-sum recovery rises from exactly $0.000$ when every process term is zeroed to $1.000$ under the full formula, and mean $\tau$ rises in parallel from $0.744$ to $0.833$.

\paragraph{Sanity check.}
The full-formula variant (all seven terms active) reproduces production exactly: $\tau = 0.833$ (set by the saturation floor on the injection set, not by ablation), recovery $= 1.000$, and mean $\lvert \Delta \textnormal{magnitude} \rvert = 0.000$ on every session. Comparison and production use the same system.

\paragraph{Monotone by active-term count.}
Table~\ref{tab:results-ablation-monotone} groups the $5{,}760$ rows by how many process-score terms are active. Both $\tau$ and sum-over-sum recovery increase with each added term, with no jumps. The key sanity check is the all-stripped row: when every process term is zeroed, total magnitude across all $45$ sessions drops cleanly to zero, so there is no baseline magnitude independent of the weights and every term does real work.

\begin{table}[htbp]
\centering
\small
\begin{tabular}{rrrr}
\hline
\textbf{active terms} & \textbf{mean $\tau$} & \textbf{sum/sum} & \textbf{n} \\
\hline
0 (all stripped) & 0.744 & \textbf{0.000} & 45    \\
1                & 0.743 & 0.170 & 315   \\
2                & 0.746 & 0.327 & 945   \\
3                & 0.754 & 0.473 & 1{,}575 \\
4                & 0.765 & 0.612 & 1{,}575 \\
5                & 0.782 & 0.745 & 945   \\
6                & 0.805 & 0.874 & 315   \\
7 (full)         & 0.833 & \textbf{1.000} & 45    \\
\hline
\end{tabular}
\caption{Mean $\tau$ and sum-over-sum magnitude recovery grouped by how many process-score terms are active. Both statistics increase monotonically with the number of active terms; sum/sum reaches $1.000$ under the full seven-term formula and drops cleanly to $0.000$ when every term is zeroed.}
\label{tab:results-ablation-monotone}
\end{table}

\paragraph{Per-term single-knob recovery.}
Table~\ref{tab:results-ablation-solo} shows how much magnitude each term contributes when it is the only active term. The \texttt{length} and \texttt{broken} terms are the strongest single-term contributors by sum/sum at $0.328$ and $0.302$, each recovering roughly a third of the total magnitude on its own; \texttt{manualIde} sits next at $0.261$. The lag term lands fourth at $0.132$, confirming the methodology's prediction that the intermediate-cleanliness penalty does measurable work on its own without dominating any existing contributor. \texttt{skipTests} and \texttt{commitGap} are mid-table at $0.089$ and $0.080$ respectively, and the \texttt{gain} term recovers $0.000$ on its own despite having the largest production weight ($50$). The latter is the regulariser signature: $\Delta C$ is comparable across user and alternative on most sessions, so \texttt{gain}-alone contributes no rank-shaping magnitude when the penalty terms it normally offsets are absent. The gap between \texttt{commitGap} ($0.080$) and \texttt{length} ($0.328$) shows recovery depends on both the weight and how much the signal varies across sessions, not on weight alone.

\begin{table}[htbp]
\centering
\small
\begin{tabular}{lrrr}
\hline
\textbf{term} & \textbf{mean $\tau$} & \textbf{mean recovery} & \textbf{sum/sum} \\
\hline
\texttt{length}    & 0.744 & 0.578 & \textbf{0.328} \\
\texttt{broken}    & 0.789 & 0.553 & 0.302          \\
\texttt{manualIde} & 0.744 & 0.442 & 0.261          \\
\texttt{lag}       & 0.669 & 0.423 & 0.132          \\
\texttt{skipTests} & 0.718 & 0.398 & 0.089          \\
\texttt{commitGap} & 0.796 & 0.368 & 0.080          \\
\texttt{gain}      & 0.744 & 0.311 & \textbf{0.000} \\
\hline
\end{tabular}
\caption{Per-term single-knob recovery on the injection set (each term active alone, the other six zeroed). The \texttt{length} and \texttt{broken} terms are the strongest single-term contributors; the lag term contributes $0.132$ on its own, fourth-largest of the seven. The \texttt{gain} term recovers $0.000$ alone because the cleanliness delta is comparable across user and alternative on most sessions.}
\label{tab:results-ablation-solo}
\end{table}

\paragraph{Per-term leave-one-out recovery.}
Table~\ref{tab:results-ablation-loo} shows what happens when you remove one term from the full seven-term formula. Removing \texttt{length} causes the largest drop in recovery ($0.718$), so the step-savings bonus contributes more single-term magnitude than any other term. Removing \texttt{gain} pushes recovery above $1.000$ to $1.060$ - the most surprising result: without \texttt{gain}, total magnitude is larger than at production. The unit being summed here is the per-divergence-point magnitude defined in \S\ref{sec:methodology-score}, i.e.\ the trajectory-final $J(\tau^\ast) - J(\tau)$ delta between each alternative and the user it is compared against; the sum is over $\lvert \text{magnitude} \rvert$ across every surfaced divergence point, including the ones whose alternative scored worse than the user (negative magnitudes contribute their absolute value). The mechanism behind the inflation is that \texttt{gain} $\cdot \Delta C$ is small but typically positive for alternatives, so it partially offsets the penalty terms; remove it and the penalty terms determine the alt-versus-user delta unopposed, producing larger per-DP absolute magnitudes which accumulate in the sum. Were the experiment to filter to alt-beats-user DPs only, removing \texttt{gain} would in fact lower the total, since alternatives sitting just above the beats-user threshold flip below it without the offset. The \texttt{gain} term therefore plays a regularising role against the penalty terms rather than acting as a positive contributor in isolation. Removing the lag term costs $0.158$ on sum/sum (production $1.000 \to 0.842$) while preserving most of the ranking ($\tau = 0.803$); the lag knob carries real magnitude into the divergence-point ranking without dominating which divergence points top the list on the injection set. Removing \texttt{skipTests} or \texttt{commitGap} reshuffles the ranking somewhat but moves recovery less than removing \texttt{length} or \texttt{manualIde} does, so these terms primarily affect which divergence point is biggest rather than how big.

\begin{table}[htbp]
\centering
\small
\begin{tabular}{lrrr}
\hline
\textbf{removed term} & \textbf{mean $\tau$} & \textbf{mean recovery} & \textbf{sum/sum} \\
\hline
\texttt{length}    & 0.833 & 0.817 & \textbf{0.718} \\
\texttt{manualIde} & 0.833 & 0.911 & 0.822          \\
\texttt{broken}    & 0.731 & 0.966 & 0.833          \\
\texttt{lag}       & 0.803 & 0.880 & 0.842          \\
\texttt{skipTests} & 0.815 & 0.903 & 0.899          \\
\texttt{commitGap} & 0.791 & 0.956 & 0.943          \\
\texttt{gain}      & 0.826 & 0.991 & \textbf{1.060} \\
\hline
\end{tabular}
\caption{Per-term leave-one-out recovery (six terms active, one removed). Removing \texttt{length} causes the largest sum/sum drop; removing \texttt{gain} inflates sum/sum above $1.000$, indicating that \texttt{gain} regularises against the penalty terms rather than acting as a positive contributor in isolation. Removing the lag term costs $0.158$ on sum/sum while leaving the ranking mostly intact ($\tau = 0.803$).}
\label{tab:results-ablation-loo}
\end{table}

\paragraph{Cross-set rerun on the user-study sessions.}
The ablation above runs on the $45$ injection sessions, which have few divergence points per session ($\max 5$, mean $1.5$). Running the same test on the $30$ user-study sessions, which have larger rankings (up to $13$ per session, mean $4.0$, with the no-feedback baseline arm carrying most of the high-count sessions), produces a partly different importance ordering. Table~\ref{tab:results-ablation-user-loo} compares leave-one-out $\tau$ across both sets. The dominant rank-carrier on the user-study set is \texttt{length}: removing it drops $\tau$ to $0.606$, the lowest figure in the table and a substantially larger drop than any other term on either set. \texttt{manualIde} and \texttt{commitGap} are the next two terms by user-study impact ($\tau = 0.657$ and $0.713$ on removal); \texttt{skipTests}, \texttt{gain}, and \texttt{broken} cluster between $\tau = 0.689$ and $\tau = 0.736$. The lag term has the smallest user-study impact ($\tau = 0.820$ on removal), so its role on the user-study set is closer to a magnitude contributor than a rank shaper, in contrast to the injection set where its leave-one-out $\tau$ is $0.803$ - comparable in magnitude to the other terms. On the injection set the figures are flatter overall ($\tau$ between $0.731$ and $0.833$), consistent with shorter per-session rankings offering less surface area for any single term to dominate. The cleanliness-gain term \texttt{gain} sits near the top of both columns ($\tau = 0.826$ on injection, $0.736$ on user-study), so it shapes the ranking more than the single-knob sweep suggested.

\begin{table}[htbp]
\centering
\small
\begin{tabular}{lrr}
\hline
\textbf{removed term} & \textbf{injection $\tau$} & \textbf{user-study $\tau$} \\
\hline
\texttt{length}    & $0.833$ & $\mathbf{0.606}$ \\
\texttt{manualIde} & $0.833$ & $0.657$ \\
\texttt{commitGap} & $0.791$ & $0.713$ \\
\texttt{skipTests} & $0.815$ & $0.719$ \\
\texttt{gain}      & $0.826$ & $0.736$ \\
\texttt{broken}    & $0.731$ & $0.689$ \\
\texttt{lag}       & $0.803$ & $0.820$ \\
\hline
\end{tabular}
\caption{Leave-one-out ranking stability against the injection sessions (Table~\ref{tab:results-ablation-loo}) and the user-study sessions, sorted by user-study impact. The \texttt{length} term is the dominant rank-carrier on the user-study set ($\tau = 0.606$ on removal). The lag term has the smallest user-study impact of the seven terms but a measurable injection-set impact ($\tau = 0.803$).}
\label{tab:results-ablation-user-loo}
\end{table}

\paragraph{Caveats.}
Five key limits apply. The cleanliness sub-weights are not ablated in this sweep because they had almost no effect on top-$1$ stability under sensitivity (\S\ref{sec:results-sensitivity}); the seven ablated terms are exactly the process-side weights of \S\ref{sec:methodology-score}. Saturation has only a small effect on this experiment: the mean saturated-divergence-point count is $0.20$ to $0.29$ across variants, so saturation does not greatly reduce the signal. $\tau$-b is unstable on small divergence-point counts, and many sessions have one or two divergence points where $\tau$ is $1.0$; mean $\tau$ at low active-counts partly reflects this floor, not pure ranking quality. The \texttt{gain}-alone row in Table~\ref{tab:results-ablation-solo} ($\tau = 0.744$, sum/sum $0.000$) sits at the all-stripped baseline because the cleanliness delta is comparable across user and alt on most sessions; the regulariser role shown in the leave-one-out table explains the same effect from the other direction. Finally, the active-count $= 0$ floor lands cleanly at sum/sum $= 0.000$ because every process-side term including lag is zeroed at that point; mean $\tau = 0.744$ at that floor reflects the $\tau$-b ties on single-DP sessions rather than residual ranking signal from a held-out weight.


\subsection{Reordering as a window onto intermediate quality}
\label{sec:results-headline}

The most informative finding is what reordering tells us about \emph{when} cleanliness arrives. The reorder synthesiser tested $194$ permutations across the $45$ sessions, which collapsed into $24$ distinct reorder windows. Under the production score formula, $10$ of those $24$ windows ($41.7\%$; Table~\ref{tab:results-divergence-quality}) have a permutation that strictly beats the user's order; a further $2$ windows acquired \emph{negative} magnitude, meaning the synthesised alternatives front-loaded \emph{less} cleanliness than the user did. Both numbers depend critically on the lag term $W_{\textsc{lag}}$ in the score: with the lag term removed, only $4$ of the $24$ windows produce a positive magnitude and the negative-magnitude verdicts collapse to zero. The lag term is therefore the operative lever for the \emph{Ordering} divergence-kind, and the headline finding of this experiment.

The mechanism is straightforward. Most IDE refactoring pairs are designed to be independent~\cite{mens2007graph}: they commute, and every permutation reaches the same final state with the same final cleanliness metrics. Under an endpoint-only score, the gain term $W_{\textsc{g}}\,\Delta C$ is identical across permutations and ordering becomes invisible. The lag term changes this: it measures the per-checkpoint deficit $\max(0, C(s_T) - C(s_i))$, so a permutation that front-loads the cleanliness-raising step accumulates a smaller running mean than one that defers it. Two permutations with the same endpoints can now score differently if their intermediate cleanliness trajectories differ -- which is precisely the discrimination an endpoint-only formula was blind to. The negative-magnitude verdicts on $2$ of the $24$ windows are the same mechanism running in reverse: on those windows the user front-loaded cleanliness better than the synthesised alts, and the lag term correctly reports the user's order as cheaper.

This has theoretical support that cuts both ways. Mens, Taentzer and Runge~\cite{mens2007graph} show that refactorings as graph transformations with parallel independence (empty critical pairs) reach the same final state under every permutation, and IDE refactorings are designed to have this property -- which explains why endpoint-only formulations score them identically. Prior empirical work either claims order matters based on interactions between code smells~\cite{liu2009smell} or tests on a single small codebase without checking AST equivalence~\cite{khrishe2016empirical}. Our $45$-session, $194$-permutation result is, to our knowledge, the first systematic evidence on real IDE sessions that independent refactorings differ measurably in process quality even when they converge on the same final code -- as long as the score formula has an intermediate-quality term.

Three things follow. First, ordering of independent refactorings matters more than endpoint-only formulations suggested, but only when the score notices intermediate quality; the $4 \to 10$ shift in positive-magnitude windows when $W_{\textsc{lag}}$ is included is the direct empirical signal. Second, the synthesiser still finds $7$ windows of non-commuting pairs (e.g.\ inline-then-rename versus rename-then-inline) where the final state itself differs; these account for most of the remaining magnitude. Third, the chapter's central methodological argument -- that process-quality evaluation cannot be reduced to endpoint comparison -- is supported here in the strongest way the data permits: a single weight added to the formula moves a controlled measurable quantity (the count of strictly-beats-user reorder windows) from $4/24$ to $10/24$, with the rest of the formula untouched.


\section{Testing the divergence-point detector}
\label{sec:results-detector}

This section evaluates the divergence-point detector against external ground truth. Precision and recall numbers on detector output are only meaningful if the per-kind labels they are computed against are themselves defensible. \S\ref{sec:results-kappa} therefore reports the inter-rater agreement results for the labelling protocol described in \S\ref{sec:methodology-rater}. \S\ref{sec:results-divergence} then measures the detector's per-kind precision and recall against those labels on the 45 injection sessions. \S\ref{sec:results-scoped-out} discusses one known recall weakness in the resulting numbers that is deliberately left open.

\subsection{Inter-rater reliability}
\label{sec:results-kappa}
Table~\ref{tab:results-userstudy-kappa} reports per-kind Cohen's $\kappa$ for the three rater pairs on the 45-session set, using the labelling protocol and interpretation thresholds defined in \S\ref{sec:methodology-rater}. All four divergence kinds reach at least substantial agreement. Ordering and IDE-replay have $\kappa = 1.00$ across every pair, meaning all three labellers agree on every session. Rework has $\kappa = 0.86$ between us and each independent rater, and $\kappa = 1.00$ between P1 and P2. Hygiene ranges from $\kappa = 0.72$ to $\kappa = 0.86$, depending on the rater pair.

\begin{table}[htbp]
\centering
\small
\begin{tabular}{lrrr}
\hline
\textbf{kind} & \textbf{Author vs P1} & \textbf{Author vs P2} & \textbf{P1 vs P2} \\
\hline
\emph{Ordering}    & $1.000$ & $1.000$ & $1.000$ \\
\emph{IDE-Replay} & $1.000$ & $1.000$ & $1.000$ \\
\emph{Rework}      & $0.861$ & $0.861$ & $1.000$ \\
\emph{Hygiene}     & $0.848$ & $0.723$ & $0.862$ \\
\hline
\end{tabular}
\caption{Per-kind Cohen's $\kappa$ on the 45-session set between us and the two independent raters and between the two independent raters. Ordering and IDE-replay are perfect on every pair; rework and hygiene reach substantial-to-almost-perfect agreement on every pair under Landis and Koch's interpretation thresholds.}
\label{tab:results-userstudy-kappa}
\end{table}

The two imperfect kinds show specific disagreement patterns. Rework disagreements occur on exactly two sessions out of $45$: session $024$ (borderline suboptimal-ordering) has a rework label from both raters but not from us, while session $043$ (borderline add-then-revert) has a rework label from us but neither rater. Since P1 and P2 agree with each other on both, our labels on these two sessions may be the outliers. Hygiene disagreements cluster on three manual-move-method sessions ($014$, $015$, $016$), where P2 added a hygiene label that neither we nor P1 added, and on session $039$ (borderline no-commit-stretch), where we added hygiene but neither rater did. Following the protocol in \S\ref{sec:methodology-rater}, these disagreements are reported as findings rather than resolved by editing labels after the fact.


\subsection{Divergence experiment}
\label{sec:results-divergence}

The third experiment evaluates the divergence-point detector itself against ground truth, using the protocol defined in \S\ref{sec:methodology-detector-eval}. The full sweep over the $45$ recorded sessions takes roughly $5$ seconds of wall-clock time. The headline numbers, expanded in the parts below, are: perfect precision on all four kinds (rework, hygiene, ordering, and IDE-replay all at $1.00$); $26$ of $45$ injections caught with a positive-magnitude alternative, every one of them ranked at top-$1$ within its own session; and -- the most informative single signal -- the count of \emph{Ordering} divergence-points whose synthesised alternative strictly beats the user's trajectory rises from $4$ to $10$ once the lag term $W_{\textsc{lag}}$ is in the formula, while the corresponding counts for the other three kinds are unchanged. The lag term is therefore the operative lever for the \emph{Ordering} kind: it does not manufacture new divergence-points (the total count of $24$ is the same with or without it), it turns previously-flat (magnitude $= 0$) ordering alternatives into recognisably-better ones.

\paragraph{Precision and recall.}
Table~\ref{tab:results-divergence-prec-recall} reports per-kind precision and recall. Rework and hygiene are perfect on both axes: $9$ of $9$ rework injections and $13$ of $13$ hygiene injections are caught at $1.00$ precision. IDE-replay is also precision-perfect. Ordering has perfect precision but only $0.36$ recall on any-expected (and $0.40$ recall on injection-only), because the reorder synthesiser's \texttt{splitOnInvalid} validator check (\S\ref{sec:methodology-reorder}) disqualifies any window containing a step the synthesiser engine cannot reproduce, and on this session set most ordering-injection windows contain at least one such step. The recall gap is documented and deliberate - its rationale is in \S\ref{sec:results-scoped-out}. Per-kind inter-rater reliability of these labels, against two independent raters, is reported in \S\ref{sec:results-kappa}.

\begin{table}[htbp]
\centering
\small
\begin{tabular}{lrrr}
\hline
\textbf{kind} & \textbf{precision} & \textbf{recall (injection)} & \textbf{recall (any expected)} \\
\hline
\emph{Ordering}   & 1.00 & 0.40 & 0.36 \\
\emph{IDE-Replay} & 1.00 & 0.76 & 0.76 \\
\emph{Rework}     & 1.00 & 1.00 & 1.00 \\
\emph{Hygiene}    & 1.00 & 1.00 & 1.00 \\
\hline
\end{tabular}
\caption{Per-kind precision and recall on the 45-session set. The ``injection'' recall column counts only the kind the playbook explicitly injected; ``any expected'' counts any kind the session's \texttt{expected\_kinds} label includes.}
\label{tab:results-divergence-prec-recall}
\end{table}

\paragraph{Detection quality.}
A precision/recall number says nothing about whether the alternatives the detector surfaces are actually useful improvements over the user's own trajectory. Table~\ref{tab:results-divergence-quality} reports the number of divergence points the detector fires per kind, the number of alternatives it enumerates, and the fraction of fires whose best alternative beats the user's trajectory under the production score formula. Across all kinds, best-alt magnitudes range from $-2$ to $+23$ with mean $+6.07$, so the percentages in the ``beats user'' column below can be read against an absolute scale of how much each winning alternative actually improved on the user's trajectory.

Hygiene never proposes a useless alternative: every one of its $13$ divergence points is a real improvement over the user's trajectory ($100\%$). IDE-replay proposes a real improvement in $12$ of $16$ fires ($75.0\%$); the four misses break down as three sessions where the user's process score was already pinned at the lower end of the $[0,100]$ range so neither the user nor the IDE-driven alternative can score below it (the magnitude is mechanically zero rather than evidence about the \texttt{manualIde} penalty), and one session where the IDE alt scored $5$ points lower than the user's manual edit. The four misses are a property of the score-range floor, not a calibration finding about \texttt{manualIde}.

Rework beats the user in $6$ of $13$ fires ($46.2\%$); one of the seven non-beats is a magnitude-zero no-op (stripping the rework loop did not move the final score) and the other six produce an alternative with negative magnitude. Two readings sit on top of each other here. The first is implementation: the rework synthesiser is intentionally simple, stripping out the user's add-then-undo edits without modelling what the intermediate states were for - if the user added a guard, ran the tests, decided the guard was unnecessary and removed it, the synthesiser treats the whole excursion as wasted motion, but the user's trajectory may have legitimately benefited from the validation step (one more checkpoint, one fewer commit-gap, a test run that pushed cleanliness up). The second is methodological: undoing work is not inherently a bad behaviour. A careful user who paths through a slightly longer trajectory to verify intermediate correctness can, and on these sessions sometimes does, score higher than a trajectory that takes the shortest line to the same final state. The $46.2\%$ beat rate is therefore not evidence that the rework detector is miscalibrated, but evidence that the score formula correctly distinguishes ``rework that wasted process signal'' (the six beats) from ``rework that earned process signal en route'' (the seven non-beats) - and that the simple loop-stripping synthesiser cannot tell the two apart by inspection of the diff alone.

Ordering surfaces $24$ distinct divergence points from $194$ enumerated permutations, of which $10$ ($41.7\%$) have a best permutation that strictly beats the user's order and a further $2$ have a best permutation that scores \emph{worse} than the user's order under the production formula (negative magnitude; the remaining $12$ are tied at magnitude zero). The $194 / 24 / 10 / 2$ breakdown is the operational evidence behind the headline in \S\ref{sec:results-headline}: the synthesiser enumerates $194$ permutations across the session set, those collapse to $24$ reorder windows (one divergence point per from-SHA / to-SHA window, magnitude set to the best permutation in the window), and the lag term $W_{\textsc{lag}}$ is what discriminates between the windows whose intermediate cleanliness trajectories differ. Without $W_{\textsc{lag}}$, only $4$ of $24$ ordering windows would show a positive magnitude and the negative-magnitude verdicts would collapse to zero.

\begin{table}[htbp]
\centering
\small
\begin{tabular}{lrrrrr}
\hline
\textbf{kind} & \textbf{DPs} & \textbf{alts enum.} & \textbf{beats user} & \textbf{beats fraction} & \textbf{max magnitude} \\
\hline
\emph{Rework}     & 13 &  13 &  6 & \textbf{46.2\%} &  8 points \\
\emph{Hygiene}    & 13 &  13 & 13 & \textbf{100\%}  &  9 points \\
\emph{IDE-Replay} & 16 &  16 & 12 & \textbf{75.0\%} & 23 points \\
\emph{Ordering}   & 24 & 194 & 10 & \textbf{41.7\%} &  9 points \\
\hline
\end{tabular}
\caption{Per-kind detection quality. ``DPs'' is the number of divergence points the detector fires; ``alts enum.'' is the total number of alternative paths enumerated across those DPs (a single ordering DP enumerates up to $k!$ permutations of a $k$-step window); ``beats user'' is the number of DPs whose best alternative trajectory beats the user's trajectory. The \emph{Ordering} beats count rises from $4$ to $10$ once the lag term $W_{\textsc{lag}}$ is in the formula; without $W_{\textsc{lag}}$ the count collapses to $4$.}
\label{tab:results-divergence-quality}
\end{table}

\paragraph{Two-way ordering discrimination.}
Two of the $24$ ordering windows now carry negative magnitude under the production formula; session $043$ is the cleanest example. The synthesised alternative converges to the same final state as the user but front-loads less cleanliness en route, so its lag term $W_{\textsc{lag}}\,\ell(\tau)$ accumulates more penalty than the user's. Under the pre-lag formula the two trajectories scored identically (same endpoint, same gain term), and the divergence-point magnitude was structurally zero. The new formula correctly reports the user's order as the cheaper one. This is the same mechanism that drives the $4 \to 10$ positive-magnitude lift, running in reverse: $W_{\textsc{lag}}$ enables two-way discrimination on intermediate quality, not just one-way inflation of positive magnitudes.

\paragraph{Injection prominence.}
The most user-facing claim in this chapter is in Table~\ref{tab:results-divergence-prominence}: of the $26$ injections the detector caught with a positive-magnitude alternative ($8$ hygiene, $12$ IDE-replay, $5$ rework, $1$ ordering), every single one was the highest-magnitude divergence point in its session. The mean rank is $1.00$ across all four kinds. A user who only opens the top recommendation in the dashboard of \S\ref{sec:arch-dashboard} therefore always sees the deliberately bad behaviour the playbook injected, and no useful alternative of another kind masks the injection. The single ordering catch is session $022$, whose ordering-injection magnitude flips from $0$ to $+1$ once the lag term is in the formula -- the same lever that drives the $4 \to 10$ shift in positive-magnitude ordering windows reported in \S\ref{sec:results-headline}. The remaining four ordering injections still surface with magnitude zero or below, consistent with the synthesiser's reach on those windows being limited by the \texttt{splitOnInvalid} validator check rather than by score-formula insensitivity (\S\ref{sec:results-scoped-out}). The prominence number is partly a property of the controlled-injection design, since the playbook injects one dominant bad behaviour per session and that behaviour is structurally likely to be the largest divergence point; the prominence number on uncontrolled real-world sessions would almost certainly be lower. The participants' sessions described in \S\ref{sec:results-userstudy} deliberately omit pre-declared expected kinds, so no prominence number is reported for them; the per-session per-kind distribution reported there is the closest external check on the prominence claim.

\begin{table}[htbp]
\centering
\small
\begin{tabular}{lrrr}
\hline
\textbf{injected kind} & \textbf{n caught} & \textbf{@ rank 1} & \textbf{mean rank} \\
\hline
\emph{Hygiene}     &  8 &  8 & 1.00 \\
\emph{IDE-Replay} & 12 & 12 & 1.00 \\
\emph{Ordering}    &  1 &  1 & 1.00 \\
\emph{Rework}      &  5 &  5 & 1.00 \\
\hline
\end{tabular}
\caption{Injection prominence: when an injection's kind is caught, where does the corresponding divergence point sit in the session's per-magnitude ranking? Top-$1$ on every catch across all four kinds. The single ordering catch is session $022$, surfaced once $W_{\textsc{lag}}$ entered the formula.}
\label{tab:results-divergence-prominence}
\end{table}

\paragraph{Caveats.}
Three limitations qualify the divergence numbers. The session set is $45$ hand-recorded sessions, so per-cell counts are small (only $5$ ordering injections, for example); the per-kind ratios are directionally correct but not enough to support tight confidence intervals, which is why raw counts are reported alongside percentages throughout. Hygiene commit-gap divergence points have a discrete magnitude exactly equal to $W_{\textsc{cg}}$ because the commit-flip alt drops the gap count by one at trajectory-final. The $194$ enumeration count for ordering is the synthesiser's reach across the session set rather than the user-facing number, which is the $24$ surfaced divergence points.


\subsection{Scoped-out fix: ordering recall}
\label{sec:results-scoped-out}

One known weakness in the divergence numbers is deliberately left open. Ordering recall sits at $0.36$ on any-expected (equivalently $0.40$ on injection-only). The cause is the \texttt{splitOnInvalid} validator check of \S\ref{sec:methodology-reorder}: any reorder window containing a step the synthesiser engine cannot reproduce (one of the three non-valid verdicts defined there) is collapsed into singletons and contributes no permutations. The gap is left open because \S\ref{sec:results-headline} has already shown that the score formula is insensitive to commuting reorders. If the validator check were relaxed, more reorder windows would be admitted, but every newly-admitted window would (by the same commuting-refactor argument) carry magnitude zero. Recall would rise, the fraction of windows that beat the user would fall by the same amount, and the substantive finding of \S\ref{sec:results-headline} would be reinforced rather than overturned. Relaxing the check is therefore a cosmetic improvement to the recall number that strengthens the headline claim of \S\ref{sec:results-headline} rather than changing it.

\section{User and agent studies}
\label{sec:results-studies}

\S\ref{sec:results-metric} and \S\ref{sec:results-detector} test the score formula and the detector in isolation; this section tests the tool end-to-end. The end-to-end question is whether the feedback the dashboard surfaces - divergence points ranked by score-delta magnitude, with concrete alternative trajectories attached - actually changes behaviour for users who did not build the tool. \S\ref{sec:results-userstudy} reports a 30-session user study with five human participants on a new codebase, split into a with-feedback arm (three participants) and a no-feedback baseline arm (two participants). \S\ref{sec:results-userstudy-agent} reports an 8-cell, 48-session agent extension on the same codebase, included as a feasibility check on the same tooling rather than a co-equal study; the framing of that section, and the reasons the human-trained detectors do not map cleanly onto agent traces, are stated there.

\subsection{User study}
\label{sec:results-userstudy}

The 30-session user study steps outside the controlled-injection sessions to three questions. First, on sessions recorded by independent participants on a different codebase and without the playbook's controlled-injection structure, does the detector's per-kind output remain coherent? Second, do per-session divergence-point rates evolve as participants accumulate tool feedback across a multi-session arc? Third, is the trajectory consistent with the dashboard feedback driving it, or would the same trajectory have appeared without that feedback? The third question is what the baseline arm exists to answer.

Five participants, referred to throughout as P1 through P5, were recruited as senior MEng Computing students at Imperial College London. All five performed six recorded refactoring sessions each on the same order-processing codebase (\texttt{user-study-fixture/}), structurally disjoint from the library codebase used in \S\ref{sec:results-divergence}: $30$ sessions in total. P1, P2, and P3 are the with-feedback arm: at the end of each session they viewed the dashboard's divergence points and trajectory advice for the session they had just completed before starting the next one. P4 and P5 are the no-feedback baseline arm: they performed the same six-session playbook on the same codebase, but never saw the dashboard between sessions. Every other aspect of the protocol - the playbook prompts, the codebase, the plugin instrumentation, the reset script between sessions - was identical across arms. Two of the five participants (P1 and P2) also produced an independent labelling of the 45-session injection set before performing their own user-study sessions; that labelling, together with our own, feeds the three-rater inter-rater agreement reported in \S\ref{sec:results-kappa} and is not repeated here.

Arm assignment was randomised: two of the five participants were drawn uniformly at random to form the baseline arm before any participant began their sessions. The arm sizes ($n = 3$ vs $n = 2$) are too small to support hypothesis testing, and no $p$-value is claimed in this section. What the design does admit is a direction-of-effect comparison: whether the gain-stripped process score and the per-session divergence-point counts move in the same direction in both arms, or diverge.

\paragraph{Descriptive findings on the participants' sessions.}
The playbook gives each session a what-not-how prompt naming the area of code to address but never the IntelliJ action to use, the testing or commit cadence to follow, or the divergence kinds the detector surfaces.\footnote{Quotes throughout this section are paraphrased from session-end notes taken during the post-study interview; full audio transcription was not feasible within the study budget. The source notes are retained as private study artefacts and are not reproduced here.} Because the playbook deliberately omits pre-declared \texttt{expected\_kinds}, the participants' sessions admit no precision/recall claim of the form reported in \S\ref{sec:results-divergence}; what they admit is a descriptive per-kind distribution and a descriptive per-session count, which Table~\ref{tab:results-userstudy-distribution} gives in aggregate by arm and by participant.

\begin{table}[htbp]
\centering
\small
\begin{tabular}{llrrrrrr}
\hline
\textbf{arm} & \textbf{participant} & \textbf{\emph{IDE-Replay}} & \textbf{\emph{Rework}} & \textbf{\emph{Hygiene}} & \textbf{\emph{Ordering}} & \textbf{total} & \textbf{DP / session} \\
\hline
with-feedback & P1 & $12$ & $2$ &  $5$ & $2$ & $21$ & $3.5$ \\
with-feedback & P2 &  $9$ & $0$ &  $8$ & $1$ & $18$ & $3.0$ \\
with-feedback & P3 &  $5$ & $0$ &  $5$ & $0$ & $10$ & $1.7$ \\
\hline
\textbf{with-feedback (sub)} & \textbf{$n = 18$ sessions} & $\mathbf{26}$ & $\mathbf{2}$ & $\mathbf{18}$ & $\mathbf{3}$ & $\mathbf{49}$ & $\mathbf{2.7}$ \\
\hline
baseline      & P4 & $21$ & $3$ & $16$ & $0$ & $40$ & $6.7$ \\
baseline      & P5 & $13$ & $0$ & $17$ & $2$ & $32$ & $5.3$ \\
\hline
\textbf{baseline (sub)}      & \textbf{$n = 12$ sessions} & $\mathbf{34}$ & $\mathbf{3}$ & $\mathbf{33}$ & $\mathbf{2}$ & $\mathbf{72}$ & $\mathbf{6.0}$ \\
\hline
\textbf{combined}            & \textbf{$n = 30$ sessions} & $\mathbf{60}$ & $\mathbf{5}$ & $\mathbf{51}$ & $\mathbf{5}$ & $\mathbf{121}$ & $\mathbf{4.0}$ \\
\hline
\end{tabular}
\caption{Per-participant divergence-point counts across the 30-session user study, grouped by arm. All five participants registered $0$ divergence points on session S04 (a cross-class duplication consolidation task that all five executed as straight manual edits) and on the small fraction of other sessions where the playbook prompt closed without a reorder window, an IDE-replayable manual sequence, or a hygiene flag. Full per-session breakdowns are in the reproducibility notebook (\texttt{tool/fixtures/notebooks/experiments.ipynb}); the per-participant totals above are the cleanest summary of the per-kind distribution.}
\label{tab:results-userstudy-distribution}
\end{table}

The aggregate per-kind distribution across the $121$ divergence points is IDE-replay $49.6\%$, hygiene $42.1\%$, rework $4.1\%$, ordering $4.1\%$ - broadly consistent with the kind proportions seen on the 45-session injection set in \S\ref{sec:results-divergence}, where IDE-replay and hygiene also dominated. Both arms exercise all four kinds, with rework and ordering sparse in both. The cleanest single number in the table is the per-session divergence-point rate: $2.7$ DPs per session in the with-feedback arm against $6.0$ in the baseline arm, a $2.2\times$ difference. The arm split also persists session-by-session - the trajectory table below shows that both arms start at comparable totals at S1 and diverge across the arc, rather than the baseline arm running at a higher rate uniformly. All five participants registered $0$ divergence points on session S04, the cross-class validation consolidation prompt, which asks the participant to consolidate three near-identical validation blocks into a single home (a delete-an-inline-block plus delete-a-dead-method task). All five executed it as straight manual edits with no rework, no skipped tests, no manual-where-IDE-could-replay candidate, and no reorder window, so the detector legitimately had nothing to surface. Used cautiously, this is the closest the chapter comes to an external true-negative observation: when the task can be discharged in a couple of clean edits and the participants do so, the detector stays quiet on all five of them. Both with-feedback participants from the original cohort identified IDE-replay as the most actionable kind in interview, with one noting that they ``didn't realise that IntelliJ could do all of these refactorings'' so didn't use them; both identified ordering as the least actionable, with one asking ``how can I actually apply that in the future without knowing''.

\paragraph{Behavioural trajectory across the six-session arc.}
The trajectory question is whether per-session divergence-point counts evolve as participants accumulate feedback - and, with the baseline arm in place, whether the same evolution is also present in participants who never saw the dashboard. Table~\ref{tab:results-userstudy-trajectory} aggregates divergence-point counts session-by-session within each arm. The with-feedback aggregate falls from $12$ in S1 to $2$ in S6, an $83\%$ reduction with a mid-arc peak at S2 driven by hygiene, then a near-monotonic decline. The baseline aggregate starts at $10$ in S1 - comparable to the with-feedback arm's $12$ at the same session - and ends at $15$ in S6, having risen rather than fallen. The two arms diverge from S3 onwards: the with-feedback arm drops from $11$ at S3 to $2$ at S6, while the baseline arm rises from $19$ to $15$ across the same range (with an S5 trough at $13$). The participants in the with-feedback arm self-reported the corresponding behavioural change in interview, with the original cohort describing how they ``ran tests after every change and committed more and used the provided automated tools'' and how they ``used the IDE refactoring methods more frequently''. The arm-level pattern aligns with those reports.

\begin{table}[htbp]
\centering
\small
\begin{tabular}{llrrrrrrr}
\hline
\textbf{arm} & \textbf{kind} & \textbf{S1} & \textbf{S2} & \textbf{S3} & \textbf{S4} & \textbf{S5} & \textbf{S6} & \textbf{$\Delta$} \\
\hline
with-feedback & \emph{IDE-Replay} & $8$ &  $6$ &  $6$ & $0$ & $5$ & $1$ & $-7$ \\
with-feedback & \emph{Hygiene}    & $3$ & $11$ &  $2$ & $0$ & $2$ & $0$ & $-3$ \\
with-feedback & \emph{Rework}     & $0$ &  $0$ &  $2$ & $0$ & $0$ & $0$ &  $0$ \\
with-feedback & \emph{Ordering}   & $1$ &  $0$ &  $1$ & $0$ & $0$ & $1$ &  $0$ \\
\textbf{with-feedback (sub)} & \textbf{total} & $\mathbf{12}$ & $\mathbf{17}$ & $\mathbf{11}$ & $\mathbf{0}$ & $\mathbf{7}$ & $\mathbf{2}$ & $\mathbf{-10}$ \\
\hline
baseline & \emph{IDE-Replay} & $5$ &  $3$ & $11$ & $0$ & $7$ & $8$ & $+3$ \\
baseline & \emph{Hygiene}    & $4$ & $11$ &  $7$ & $0$ & $6$ & $5$ & $+1$ \\
baseline & \emph{Rework}     & $0$ &  $1$ &  $1$ & $0$ & $0$ & $1$ & $+1$ \\
baseline & \emph{Ordering}   & $1$ &  $0$ &  $0$ & $0$ & $0$ & $1$ &  $0$ \\
\textbf{baseline (sub)} & \textbf{total} & $\mathbf{10}$ & $\mathbf{15}$ & $\mathbf{19}$ & $\mathbf{0}$ & $\mathbf{13}$ & $\mathbf{15}$ & $\mathbf{+5}$ \\
\hline
\end{tabular}
\caption{Per-kind divergence-point trajectory across the six-session arc, summed within each arm. The with-feedback aggregate drops from $12$ at S1 to $2$ at S6 ($\Delta = -10$); the baseline aggregate rises from $10$ to $15$ ($\Delta = +5$). The two arms are comparable at S1 and diverge from S3 onwards. Both arms register $0$ divergence points on S04 (a clean-consolidation prompt all participants discharged in a couple of edits). IDE-replay and hygiene carry most of the arm-level signal in both directions.}
\label{tab:results-userstudy-trajectory}
\end{table}

The two arms differ on which kinds carry the trajectory. In the with-feedback arm, IDE-replay falls from $8$ to $1$ and hygiene from $3$ to $0$ (with the S2 hygiene spike consistent with participants first seeing commit-cadence feedback at the end of S1 and applying it from S3 onwards). In the baseline arm, IDE-replay does not fall - it ends $+3$ on the arc - and hygiene ends $+1$ rather than $-3$. Rework and ordering are episodic in both arms: ordering fires rarely on the participants' sessions ($5$ positive-magnitude divergence points in total across $30$ sessions), consistent with the headline finding of \S\ref{sec:results-headline} that the score formula is insensitive to commuting reorders, and rework spikes only on the hardest sessions where contiguous undo/redo cycles drive several flags from a single sequence of mis-clicks. One with-feedback participant raised the rework behaviour as a metric-design complaint in interview, observing that rework ``picked up on undos + redos a lot (in a bad way) [-] over penalised'' on contiguous-step sequences caused by mistyped keys; the per-session detail behind the trajectory table is consistent with that account.

The six sessions are not interchangeable: S2 (magic-number renames) and S4 (cross-class consolidation) are short, mechanical sessions where divergence-point counts collapse for any participant regardless of arm, while S3 and S6 are structurally harder sessions where the opportunity for hygiene and rework flags is built into the playbook prompt. Cross-session comparisons of raw divergence-point counts within a single arc are therefore confounded by task difficulty - which is exactly why the baseline arm matters. The baseline arm rides the same task-difficulty envelope as the with-feedback arm (both arms see the same S2 hygiene rise and the same S4 quiet), so the divergence between the two arms across the arc is the signal that is left after the task-difficulty confound is held constant. The same confound applies to the agent extension in \S\ref{sec:results-userstudy-agent}.

\paragraph{Process scores under the production and gain-stripped weightings.}
The within-session process score itself can be partially separated from task difficulty by inspecting it twice: once under the production weights, and once under a weighting that sets the cleanliness-gain weight $W_g$ to zero. The cleanliness-gain term is the only term in the score formula that scales with how much the code itself improved over the session, and the ablation in \S\ref{sec:results-ablation} already showed that this term acts as a regularizer on divergence-point ranking rather than reshaping it. $W_g$ is the largest weight in the production formula by design, since the anti-gaming axiom of \S\ref{sec:methodology-score} requires the cleanliness reward to outweigh any process-side shortcut a developer could plausibly take by trading off code quality for procedural discipline; zeroing $W_g$ in a free-form deployment would invite exactly that gaming. The participants' sessions are not free-form: the playbook fixes the structural change each session must produce, so the cleanliness outcome is set by whether the participant completed the prompt rather than by anything the participant could trade off, and the gain-stripped view is a legitimate read on the process-discipline terms alone for this experiment. Zeroing $W_g$ leaves the broken-time, skipped-tests, manual-when-IDE, length-bonus, and commit-gap terms in place, so the remaining score asks ``did the user execute the process well'' rather than ``did the code end up cleaner''. Table~\ref{tab:results-userstudy-scores} reports the gain-stripped score across the six-session arc for all five participants, grouped by arm.

\begin{table}[htbp]
\centering
\small
\begin{tabular}{llrrrrrrrr}
\hline
\textbf{arm} & \textbf{participant} & \textbf{S1} & \textbf{S2} & \textbf{S3} & \textbf{S4} & \textbf{S5} & \textbf{S6} & \textbf{$\Delta$} & \textbf{slope} \\
\hline
with-feedback & P1 & $25$ & $22$ & $41$ & $38$ & $34$ & $48$ & $+23$ & $+4.6$ \\
with-feedback & P2 & $16$ & $13$ & $29$ & $39$ & $37$ & $47$ & $+31$ & $+6.2$ \\
with-feedback & P3 & $39$ & $15$ & $42$ & $42$ & $31$ & $48$ &  $+9$ & $+1.8$ \\
\textbf{with-feedback (mean)} & ($n = 3$) & $\mathbf{26.7}$ & $\mathbf{16.7}$ & $\mathbf{37.3}$ & $\mathbf{39.7}$ & $\mathbf{34.0}$ & $\mathbf{47.7}$ & $\mathbf{+21.0}$ & $\mathbf{+4.2}$ \\
\hline
baseline      & P4 & $22$ & $17$ &  $8$ & $22$ & $12$ & $15$ &  $-7$ & $-1.4$ \\
baseline      & P5 & $23$ & $17$ &  $9$ & $36$ & $11$ & $20$ &  $-3$ & $-0.6$ \\
\textbf{baseline (mean)}      & ($n = 2$) & $\mathbf{22.5}$ & $\mathbf{17.0}$ & $\mathbf{8.5}$ & $\mathbf{29.0}$ & $\mathbf{11.5}$ & $\mathbf{17.5}$ & $\mathbf{-5.0}$ & $\mathbf{-1.0}$ \\
\hline
\end{tabular}
\caption{Trajectory-final process score across the six-session arc for each participant under the cleanliness-gain-zeroed weighting ($W_g = 0$), grouped by arm. The production-weighted view is dominated by the cleanliness-gain term and is reported per-session in the reproducibility notebook. $\Delta$ is $S6 - S1$ and slope is $\Delta / 5$ in score units per session. The arm-mean rows aggregate across participants within the arm.}
\label{tab:results-userstudy-scores}
\end{table}

Two patterns sit in the table. First, all three with-feedback participants finish above where they started ($+23$, $+31$, $+9$), while both baseline participants finish below where they started ($-7$, $-3$). The arm means separate cleanly: $+21.0$ vs $-5.0$, or $+4.2$ per session vs $-1.0$ per session. With per-participant variance comparable to the between-arm difference, the finding is reported as a direction of effect rather than a hypothesis test - which is what $n = 3$ vs $n = 2$ admits.

Second, the within-arc shape carries the same task-difficulty signature in both arms: S2 dips in four of five participants, S4 sits near a per-participant ceiling for everyone, and the cross-arm separation builds from S3 onwards. S2 (magic-number renames executed under intermediate-cleanliness pressure) and S3 (within-class duplication consolidation under stricter test-and-commit cadence) are the two sessions where the dashboard's hygiene and IDE-replay advice would have been most actionable for a participant who saw it after S1, which is consistent with the with-feedback arm's S3 mean ($37.3$) sitting above the baseline arm's ($8.5$) by the largest single-session margin on the arc. The S4 convergence ($39.7$ vs $29.0$) reflects that S4 is a clean-execution session for everybody and offers little room for either arm to differentiate.

The playbook orders the six sessions in roughly increasing structural complexity (S1 renames through to S6 multi-step reshape). The with-feedback arm's improvement is therefore happening despite progressively harder tasks, not riding on easier ones - and the baseline arm's decline is happening over the same difficulty curve, which removes increasing difficulty as a confound for the arm-level direction. P3's lower slope ($+1.8$) is consistent with starting at a higher S1 ($39$) closer to ceiling than P1 or P2: P3 has less room to improve, not a different arc shape.

\paragraph{With-feedback vs no-feedback baseline arm.}
The arm-level comparison is the strongest causal-direction signal in the chapter, because arm assignment was randomised over a single participant pool and every other aspect of the protocol was held constant. The baseline arm rides the same task difficulty, the same codebase, the same playbook order, and the same plugin instrumentation as the with-feedback arm; the only manipulated variable between arms is whether the participant saw the dashboard between sessions. The arm-mean rows of Table~\ref{tab:results-userstudy-scores} report a $+21.0$ vs $-5.0$ gain-stripped $\Delta$ across the arc (a $26.0$-point gap, or $+4.2$ vs $-1.0$ per session), and the arm rows of Table~\ref{tab:results-userstudy-trajectory} report a $-10$ vs $+5$ divergence-point $\Delta$ across the arc (a $15$-DP gap). Both signals point the same direction.

The gain-stripped subscore is largely insulated from how familiar the participant is with the codebase itself - the terms active under $W_g = 0$ are commit cadence, skipped tests, and manual-when-IDE rate, none of which depend on knowing the code better. What the design does not separate is the with-feedback participants intuiting what the dashboard rewards and adjusting behaviour to satisfy it (a score-gaming reading), or becoming more comfortable with the playbook idiom across the arc (a task-style-familiarity reading). The threats-to-validity chapter (\S\ref{sec:threats-per-experiment}) discusses both. The baseline arm rules out one strong alternative reading - that the with-feedback arc reflects the playbook getting easier as the participant moves through it - because the baseline participants worked the same playbook in the same order and trended downward across the arc rather than upward. The remaining limit is sample size. With $n = 3$ vs $n = 2$ the arm means are estimated from too few participants to support a hypothesis test, and per-participant variance is comparable to the between-arm difference, so the result is reported as a randomised but underpowered direction-of-effect claim: ``under random arm assignment, the with-feedback arm improved across the arc, the baseline arm did not, and the within-arm variance is smaller than the between-arm difference.'' No $p$-value or confidence interval is claimed.

One additional qualitative finding does not show in the trajectory tables but matters for future tuning: hygiene also produces a false-positive shape, firing when a participant is thinking rather than stuck. One with-feedback participant flagged this in the interview as being ``penalised for not running tests when wasn't quite finished but just was thinking'', and the pattern is reflected in the S2 hygiene peak in both arms. Both this finding and the rework over-penalisation are reported here as candidate metric-tuning directions rather than detector bugs: the divergence points are accurate readings of what the trajectory contains, but the metric design treats some legitimate-pause and mistyped-key cases more harshly than the participant would.

\paragraph{Per-step inspection as a usability win.}
\label{sec:results-userstudy-stepui}
Both participants singled out the dashboard's per-checkpoint detail panel as a positive part of the interface, independently of the four detector kinds. One participant said the per-checkpoint view made it easy to ``see what action you made to cause the issues to happen'', and that the trajectory graph showing the highest score reachable from each point onwards was ``egging you on to be a better person'' (i.e.\ the alternative-trajectory overlays were read as constructive rather than punitive). The other participant said they could ``see the specific point you went wrong'', valued the cleanliness-signal ``breakdowns'' on each step, and called the dashboard ``easy to track session at specific points and code smells''. These observations validate the per-step smell-attribution and metric-shift breakdown described in \S\ref{sec:arch-metrics} as a usability feature the participants engaged with, not just an architectural by-product of how the report is assembled.


\subsection{Agent extension experiment}
\label{sec:results-userstudy-agent}

\paragraph{Framing.}
This section reports an extension experiment, not a co-equal evaluation. The tool described in this thesis is designed against human edit patterns: it captures IDE refactor-menu invocations, debounces edit bursts on a 2-second idle, and thresholds commit cadence against the 60-second composite window of \cite{murphyhill2009howrefactor}. None of these signals map cleanly onto how a coding agent emits file changes (\S\ref{sec:results-userstudy-agent-limits} below itemises the three structural biases that follow). The intent of running the playbook against eight agent stacks is therefore narrower than for the user study: to observe how the human-trained detectors behave on a structurally different actor, and to report what the corpus admits at face value without overclaiming the tool as an agent-evaluation framework.

\paragraph{Setup.}
Eight agent stacks ran the same six-session playbook used by the human participants, on the same \texttt{user-study-fixture/} codebase: $48$ sessions in total. The eight stacks combine three model families (Claude, OpenAI GPT, Gemini) with three command-line harnesses (Claude Code, OpenCode, Cursor CLI), and include two no-feedback baselines, matching the human study's arm structure. Six stacks are the with-feedback arm: between sessions, the human facilitator pasted the markdown output of the \texttt{feedbackForAgent} task (\texttt{tool/analysis/.../experiment/FeedbackForAgent.kt}) into the agent's prompt at the start of the next session. That task renders an analysis report's divergence points and trajectory advice into the same plain-text summary that a human participant would read on the dashboard, so the agent received the same information the human participants did, just in textual rather than rendered form. Two stacks are the no-feedback baseline arm (\texttt{claude-4.7-opus-medium-claude-code-baseline} and \texttt{gpt-5.5-medium-opencode-baseline}): they ran the same six-session playbook on the same codebase but were never shown the dashboard between sessions. Every other aspect of the protocol matched the with-feedback agent arm. Each agent honoured a small setup contract: tests were invoked via a wrapper script that brackets the gradle call with \texttt{TEST\_RUN\_STARTED} / \texttt{TEST\_RUN\_FINISHED} events in the plugin trace, and commits used standard \texttt{git commit} which the plugin's reflog watcher captured as \texttt{GIT\_COMMIT} events. Apart from starting and stopping the recording at session boundaries and pasting the previous session's feedback into the next prompt where applicable, no human intervened during a session.

\paragraph{Per-cell descriptive arc.}
Table~\ref{tab:results-userstudy-agent-trajectory} reports the gain-stripped trajectory-final process score across the six-session arc for each of the eight cells, along with the per-cell divergence-point total. The gain-stripped view is used here for the same reason as the human user-study table: it isolates the process-discipline terms from the cleanliness-gain term, which is dominated by playbook completion rather than process quality.

\begin{table}[htbp]
\centering
\small
\begin{tabular}{rllrrrrrrrr}
\hline
\textbf{\#} & \textbf{cell} & \textbf{arm} & \textbf{S1} & \textbf{S2} & \textbf{S3} & \textbf{S4} & \textbf{S5} & \textbf{S6} & \textbf{$\Delta$J} & \textbf{$\sum$DP} \\
\hline
1 & claude-2-opus-4.7-claude-code              & feedback & $35$ & $39$ & $40$ & $45$ & $40$ & $35$ &   $0$ &  $5$ \\
2 & claude-opus-4.7-medium-opencode            & feedback & $37$ & $36$ & $40$ & $45$ & $37$ & $38$ &  $+1$ &  $4$ \\
3 & claude-opus-4.7-thinking-medium-cursor-cli & feedback & $35$ & $39$ & $40$ & $45$ & $39$ & $39$ &  $+4$ &  $6$ \\
4 & gemini-3.5-flash-medium-opencode           & feedback & $33$ & $37$ & $40$ & $45$ & $39$ & $37$ &  $+4$ &  $9$ \\
5 & gpt-5.5-medium-cursor-cli                  & feedback & $23$ & $38$ & $39$ & $45$ & $37$ & $37$ & $+14$ &  $7$ \\
6 & openai-gpt-5.5-medium-opencode             & feedback & $35$ & $38$ & $40$ & $45$ & $36$ & $36$ &  $+1$ &  $5$ \\
\hline
7 & claude-4.7-opus-medium-claude-code-baseline & baseline & $36$ & $34$ & $39$ & $45$ & $36$ & $40$ &  $+4$ &  $3$ \\
8 & gpt-5.5-medium-opencode-baseline            & baseline & $25$ & $34$ & $40$ & $45$ & $36$ & $38$ & $+13$ &  $4$ \\
\hline
\multicolumn{3}{l}{\textbf{with-feedback (mean, $n = 6$)}} & $\mathbf{33.0}$ & $\mathbf{37.8}$ & $\mathbf{39.8}$ & $\mathbf{45.0}$ & $\mathbf{38.0}$ & $\mathbf{37.0}$ & $\mathbf{+4.0}$ & $\mathbf{6.0}$ \\
\multicolumn{3}{l}{\textbf{baseline (mean, $n = 2$)}}      & $\mathbf{30.5}$ & $\mathbf{34.0}$ & $\mathbf{39.5}$ & $\mathbf{45.0}$ & $\mathbf{36.0}$ & $\mathbf{39.0}$ & $\mathbf{+8.5}$ & $\mathbf{3.5}$ \\
\hline
\end{tabular}
\caption{Per-cell gain-stripped trajectory and total divergence-point count across the eight agent stacks. Cells 1--6 are with-feedback (named for the model $\times$ harness combination); cells 7 and 8 are no-feedback baselines, named with the \texttt{-baseline} suffix. The arm-mean rows aggregate across cells within the arm.}
\label{tab:results-userstudy-agent-trajectory}
\end{table}

Four observations follow. First, every agent cell sits at $J = 45$ at S4 - a uniform-across-cells corroboration of the user-study finding that S4 is genuinely a no-event session under the playbook. Second, all eight cells finish above where they started: the smallest $\Delta J$ in the table is $0$ (cell 1) and the largest is $+14$ (cell 5). Third, the divergence-point counts are small and almost entirely IDE-replay: across $48$ agent sessions the total is $42$ DPs, of which $41$ are IDE-replay, $1$ is ordering, and $0$ are rework or hygiene. Fourth, the per-cell DP rate ($0.7$--$1.5$ per session for the with-feedback arm, $0.5$--$0.7$ for baseline) is an order of magnitude lower than the human cohort's ($2.7$--$6.0$ per session); agents start at a much higher procedural-discipline floor than the human participants do.

\paragraph{With-feedback vs no-feedback baseline arm.}
The arm-level comparison in Table~\ref{tab:results-userstudy-agent-trajectory} does not show the same direction-of-effect signal that the human study did. The with-feedback arm gains $+4.0$ on $\Delta J$ across six cells; the baseline arm gains $+8.5$ across two cells. Both arms' slopes are positive, but the baseline arm's slope ($+1.7$ per session) is larger than the with-feedback arm's ($+0.8$). Per-session DP totals tell the same story: the with-feedback arm sums to $36$ DPs across $36$ sessions ($1.0$ per session), the baseline arm to $7$ DPs across $12$ sessions ($0.6$ per session). On this corpus the tool's between-session feedback shows no measurable causal effect on agent process scores, and what difference does appear runs in the opposite direction to the human result.

Several factors plausibly contribute to the null result, none of them indicating the tool's feedback is harmful to agents:
\begin{itemize}
    \item \textbf{Floor effect.} Agents start at a much higher procedural-discipline ceiling than humans (mean S1 gain-stripped $J \approx 33$ across cells vs $\approx 22.5$ for the human baseline arm, with baseline DP rates $\approx 0.5$/session for agents vs $\approx 6.0$/session for human baselines). There is materially less room for any intervention to move the score.
    \item \textbf{Detector-bias floor.} Three of the four divergence-point kinds rarely fire on agent traces at all (see \S\ref{sec:results-userstudy-agent-limits} below). Whatever between-session adjustment a with-feedback agent makes on those kinds is invisible to the score by construction.
    \item \textbf{Anomalous S1 starts in the baseline arm.} The $n = 2$ baseline arm includes \texttt{gpt-5.5-medium-opencode-baseline} starting at $S1 = 25$ (the lowest S1 in the table), which inflates that cell's $\Delta J$ to $+13$. With one of the two baseline cells starting that far below ceiling, the baseline arm mean is sensitive to a single low-S1 reading.
\end{itemize}

\paragraph{Cross-stack effects.}
The six with-feedback cells span a model axis (Claude vs Gemini vs GPT) at fixed harness and a harness axis (Claude Code vs OpenCode vs Cursor CLI) at fixed model. With one cell per axis position there is no inferential weight to give the comparison, but the descriptive shape is worth recording. On the model axis with the OpenCode harness held out, Claude (cell 2), Gemini (cell 4), and GPT (cell 6) produce slopes $+0.20 / +0.80 / +0.20$ and DP totals $4 / 9 / 5$: Gemini surfaces the most DPs but climbs $J$ the fastest, consistent with a willingness to commit mid-task partial work that surfaces as IDE-replay but recovers by S6. On the harness axis with the Claude model held out, Claude Code (cell 1), OpenCode (cell 2), and Cursor CLI (cell 3) produce slopes $0.00 / +0.20 / +0.80$ and DP totals $5 / 4 / 6$: Cursor CLI improves fastest at fixed model. The cross-stack comparison (cell 1 vs cell 6, Claude / Claude Code at slope $0.00$ vs GPT / OpenCode at slope $+0.20$, $5$ DPs each) is roughly tied.

\paragraph{Detector-bias limitations on agent traces.}
\label{sec:results-userstudy-agent-limits}
Three structural reasons explain why the human-trained detectors under-fire on agent traces:

\begin{itemize}
    \item \textbf{No IDE-driven refactorings.} Agents apply file edits via text tool calls, not via IntelliJ's Refactor menu. Every structural refactoring an agent performs is classified as a manual edit by the detector, so the \emph{IDE-Replay} kind fires on the structural refactorings that an IntelliJ-side refactor would have produced - but the agent has no actuator to switch to even if it reads the feedback. This is the dominant kind in the agent corpus ($41$ of $42$ DPs) and is the cleanest example in the chapter of feedback being correctly interpreted but structurally inactionable.
    \item \textbf{Commit-gap threshold rarely triggered.} The \emph{Hygiene} kind's commit-gap detector flags six refactor checkpoints elapsed without a commit. Agents batch many file changes per checkpoint and arrive at S6 in fewer total checkpoints than a human session of equivalent length. In a six-session arc the threshold is rarely reached, which is consistent with the $0$ hygiene DPs across all $48$ agent sessions.
    \item \textbf{Edit-burst debouncer collapses each agent turn into one burst.} The edit-burst tracker debounces on \texttt{DEBOUNCE\_MS = 2000L} (2 s; see \texttt{tool/ide-plugin/.../EditBurstTracker.kt}). It does fire on tool-applied file writes. But agents emit many file edits in a tight burst, then pause for reasoning longer than 2 s; each tool turn therefore collapses to one burst event, compressing the burst-vs-checkpoint signal that the \emph{Rework} and \emph{Ordering} kinds depend on. This is consistent with the $0$ rework DPs and $1$ ordering DP across the $48$ agent sessions.
\end{itemize}

None of these are tool bugs - the detectors do the right thing on human traces. They imply that the agent results in Table~\ref{tab:results-userstudy-agent-trajectory} should be read as descriptive feasibility evidence rather than a verdict on the tool's intended use, which is human-developer refactoring.

\paragraph{Feedback engagement varies by model family.}
The aggregate arm comparison says feedback did not move the score. Per-cell transcripts say something more interesting. Skimming the \texttt{transcript.md} file at the root of each with-feedback cell directory shows a sharp split. The three Claude-family cells (1, 2, 3) explicitly engaged with the feedback in writing. Cell 1 is the strongest case: after S1 warnings about \texttt{BUILD\_OFTEN\_BROKEN}, \texttt{TEST\_REGRESSIONS}, and \texttt{PROCESS\_SCORE\_DEGRADED}, S2 opens with ``commit smaller and more often - one commit per rename (or per logical step), not one at the end. Test after each step.'' After S3 flagged \texttt{REFACTORING\_INTRODUCED\_SMELLS}, S4 planning included ``when extracting a method, leave the extracted body cleaner - extract and tidy, not extract and stop.'' Cell 2 (\texttt{claude-opus-4.7-medium-opencode}) follows the same pattern: after S1, ``\texttt{BUILD\_OFTEN\_BROKEN} - likely caused by intermediate states - I'll land each rename atomically per checkpoint.'' Cell 3 (\texttt{claude-opus-4.7-thinking-medium-cursor-cli}) shows quieter engagement but visible adjustment after the S5 IDE-replay summary. The three non-Claude cells (4 Gemini, 5 GPT-5.5 Cursor CLI, 6 GPT-5.5 OpenCode) show no visible engagement: they proceed with consistent (and competent) behaviour session-to-session but do not acknowledge or react to the between-session feedback in their conversation traces.

The split is descriptively interesting on its own, and is a reasonable starting point for a follow-up study on how different model families treat reflective critique in iterative tasks. It does not, however, rescue a causal claim about the tool's feedback driving score improvement. Even on the Claude-family cells, where the transcripts show the feedback being read and acted on, the gain-stripped $\Delta J$ does not separate from the baseline arm.

\paragraph{Qualitative threads from cell 1 (\texttt{claude-2-opus-4.7-claude-code}).}
The aggregate trajectory table flattens what is, on inspection, a substantively richer story for cell 1 specifically. The agent's written reflections after each session show it reading the feedback, diagnosing its own anti-patterns, and adjusting its approach between sessions in ways the per-kind counts only partially capture. These threads are reported as cell-1 anecdotes that illustrate what feedback engagement looks like in transcript form, rather than as findings that generalise across the eight-cell agent extension.

\emph{Thread 1: commit cadence, over-corrected and then resolved.} S1 produced cell 1's worst-scoring session of the arc: one bulk commit at the end of all three renames, a session-long stretch without committing, a $28\%$ build-broken rate (the build was broken on $2$ of $7$ checkpoints), one test regression, and a production-weighted process score that fell from $50$ to $0$ - the largest single-session drop on the arc. The post-session feedback fired three advice items at critical and warning levels. The agent's response, written before S2 began, was the direct procedural correction quoted above. S2 implemented exactly this: three steps as three separate \texttt{git commit} calls, each preceded by \texttt{./agent-test}; the resulting session produced $0$ divergence points, $0$ advice items, and the production-weighted process score climbed from $50$ to $79$. The agent then over-corrected in the opposite direction. Its pre-S3 statement was ``I'll aim to commit even more granularly (each individual constant/edit, not each step)'', and S3 fired three IDE-replay divergence points (one at magnitude $-6$), \texttt{REFACTORING\_INTRODUCED\_SMELLS}, and a production-weighted score drop from $50$ to $30$. By S6 the agent had revised the rule a second time, but in a more interesting direction: the S6 advice flagged an $8$-checkpoint commit stretch and a $37\%$ build-broken rate ($3$ of $8$ checkpoints), and the agent's diagnosis identified the cause as a structural tension the rule itself created: ``for an Extract Method, do it as a single atomic edit (insert new method first, then replace call site) so the build never breaks mid-step.'' The agent had discovered, by being penalised for it twice, that ``commit more often'' and ``don't break the build'' are in conflict when an edit is structurally multi-step, and resolved the conflict by switching to larger atomic diffs rather than finer-grained commits.

\emph{Thread 2: smell-introduction feedback produced a specific self-critique.} After S3 the agent fired \texttt{REFACTORING\_INTRODUCED\_SMELLS} alongside the process-score drop, and the agent's written response identified its own anti-pattern by name: ``methods like \texttt{chargeCash}/\texttt{refundVoucher} are arguably trivial wrappers that didn't earn their keep,'' and ``could have made \texttt{formatAddress} cleaner.'' The follow-on sessions support the diagnosis: S4 produced zero divergence points and zero advice items, S5 produced two IDE-replay flags but no \texttt{REFACTORING\_INTRODUCED\_SMELLS}, and the smell-introduction advice did not fire again until S6 - where the agent attributed it to multi-edit extract-method mechanics rather than the trivial-wrapper anti-pattern from S3.

\emph{Thread 3: the agent acknowledges the structural IDE-replay limitation.} The persistent failure mode for cell 1 is IDE-replay. All five of cell 1's divergence points are IDE-replay flags, and the count does not decline across the arc; it concentrates on the structurally heavier sessions (S3, S5, S6) where the agent performed enough structural refactoring to surface any divergence points at all. The agent itself acknowledges the structural nature of this in its S6 reflection, observing that ``the IDE-replay flags persist - manual extracts consistently underperform the IntelliJ refactoring; in a real run I'd use the Extract Method action.'' The signal is real and the recommendation correct, but the agent lacks the tool to execute it - which is the same general observation that \S\ref{sec:results-userstudy-agent-limits} states as a population-level structural limitation rather than a single-cell anecdote. A setup that exposed IntelliJ's refactor actions as agent tool calls would close the IDE-replay gap; that work is named in the future-work section as the main follow-up for the agent extension.

